% 第九卷 10.1.4.2 节：W2NER词对关系网格（教学句“左／和／右／肺／均／清晰”）。
% 行为i、列为j；上三角(j>i)放NNW，下三角(i>j)放THW-类型，对角线可放单词实体的THW。
% 正标签：G(1,4)=G(3,4)=NNW，G(4,1)=G(4,3)=THW-BODY，其余为NONE。
\begin{tikzpicture}[
  x=1mm,y=1mm,figure picture,
  cellbox/.style={draw=FigureGrid,line width=.4pt,minimum size=8mm,inner sep=0pt},
  upper/.style={cellbox,fill=white},
  lower/.style={cellbox,fill=FigurePaper},
  nnw/.style={cellbox,fill=FigureModelFill,draw=FigureModel,line width=.7pt,
    font=\figurefont\scriptsize},
  thw/.style={cellbox,fill=FigureInputFill,draw=FigureInput,line width=.7pt,
    font=\figurefont\tiny,align=center},
  lab/.style={font=\figurefont\footnotesize,text=FigureInk},
  note/.style={font=\figurefont\scriptsize,text=FigureMuted,anchor=west,align=left}
]
  % 网格底色
  \foreach \i in {1,...,6}{
    \foreach \j in {1,...,6}{
      \pgfmathtruncatemacro{\up}{\j>\i ? 1 : 0}
      \pgfmathtruncatemacro{\dn}{\i>\j ? 1 : 0}
      \ifnum\up=1
        \node[upper] at ({(\j-1)*8},{-(\i-1)*8}) {};
      \else\ifnum\dn=1
        \node[lower] at ({(\j-1)*8},{-(\i-1)*8}) {};
      \else
        \node[cellbox,fill=white] at ({(\j-1)*8},{-(\i-1)*8}) {};
      \fi\fi
    }
  }
  % 行列标签
  \foreach \k/\w in {1/左,2/和,3/右,4/肺,5/均,6/清晰}{
    \node[lab] at ({(\k-1)*8},7) {\w};
    \node[font=\figurefont\scriptsize,text=FigureMuted] at ({(\k-1)*8},11.5) {$j{=}\k$};
    \node[lab,anchor=east] at (-6,{-(\k-1)*8}) {\w};
    \node[font=\figurefont\scriptsize,text=FigureMuted,anchor=east] at (-13,{-(\k-1)*8}) {$i{=}\k$};
  }
  % 正标签
  \node[nnw] at ({(4-1)*8},{-(1-1)*8}) {NNW};
  \node[nnw] at ({(4-1)*8},{-(3-1)*8}) {NNW};
  \node[thw] at ({(1-1)*8},{-(4-1)*8}) {THW\\BODY};
  \node[thw] at ({(3-1)*8},{-(4-1)*8}) {THW\\BODY};
  % 区域说明
  \node[note] at (50,0) {上三角（$j>i$）：NNW，\\实体内部的次序相邻词};
  \node[note] at (50,-14) {下三角（$i>j$）：THW-类型，\\由尾词指回首词};
  \node[note] at (50,-28) {解码：THW(4,1)与路径1$\to$4\\得“左…肺”$[0,1)\cup[3,4)$；};
  \node[note] at (50,-38) {THW(4,3)与路径3$\to$4\\得“右肺”$[2,4)$};
\end{tikzpicture}
